% Replacement prose for the integrated manuscript, Fable 5.2 xhigh final review.
% Pinned target: integrated/manuscript.pdf sha256 8268234265df2bb9e0cf9a3303a51f79992e96e69a42b1d074c867697313feb5.
% Each block names its finding in findings.json, the target file and line, and what it replaces.
% Every number below was re-derived from existing files; see verification.md for the script and source.
% This file is body text only. build/replacement_build.tex wraps it for a standalone syntax check.
% Nothing here was applied to any shared file.

%======================================================================================
\ReplacementBlock{FF-D01}{sections/development.tex lines 6, 14, 16 and 18}
% Replaces the last sentence of the first paragraph, extends the second and third, and replaces
% the final sentence of the fourth ("In a received conflict, content may request a mortgage ...").

% --- line 6, last sentence ---
The twenty residual windows contain nine in which no recorded menu offers construction on the completed group, four defaults, four received-call conflicts and three other valid ends (Figure~\ref{fig:execution}).

% --- line 14, append after "Bank houses are available in all nine." ---
The category concerns the completed group, not every group. In Grok's two G1 windows the menu does offer a house on Atlantic Avenue, in a yellow group Grok already owns.

% --- line 16, replace the sentence beginning "In all four Grok windows" ---
In all four Grok windows, the attached note explicitly defers redemption or construction for stated liquidity reasons, and both G1 notes also decline the available Atlantic house.

% --- line 18, replace the final sentence with this paragraph ---
The four conflict windows are all Grok's, and they differ in what the content names. In G8, G12 and G13, Grok first mortgages other deeds within the window, raising \$310, \$350 and \$335. The content of the final decision then names a construction action, and the received call ends the turn. The G8 content lists one house on each orange deed, \$300 against \$472 in cash. The G13 content lists one house on each of two yellows, \$300 against \$357. Every deed named in those two plans is on the menu's buildable list, and the bank holds at least twenty-five houses. The G12 content lists three, two and two pink houses for \$700 against \$726, a count-valued plan that per-entry validation would reject. In G20 the content names a mortgage. No content plan was executed or checked by the engine, so these are readings of the recorded menus, not witnesses. Grok builds on orange fifteen turns later in G8 and on pink four turns later in G12. It never builds in G13.

%======================================================================================
\ReplacementBlock{FF-D01 (figure, methods and appendix wording)}{tools/build\_exhibits.py line 48; sections/development.tex line 10; sections/methods.tex line 8; sections/appendix.tex lines 94 and 97}

% Figure 1 bar label: "No option observed" becomes
No option for group

% Figure 1 caption, last sentence becomes
Absence of a construction option for the completed group on recorded menus does not rule out unvisited financing choices.

% methods.tex line 8, last sentence becomes
A menu sequence that never offers construction on the completed group does not establish impossibility under unvisited financing choices.

% appendix.tex line 94, subsection title becomes
Financing budgets where no construction option for the group appeared

% appendix.tex line 97, second sentence becomes
Every target group is complete, has no buildings, and has bank houses available; the two G1 prompts also list a buildable house on another group.

%======================================================================================
\ReplacementBlock{FF-D02}{sections/development.tex line 20}
% Replaces the second sentence of the paragraph beginning "Across all management menus offering construction".

Sixty-six are defaults and twenty-four conflict with the action named in content. In twenty-two of those twenty-four, the content names a construction action.

%======================================================================================
\ReplacementBlock{FF-L01}{sections/finance\_inheritance.tex line 18}
% Replaces the whole paragraph beginning "Actual staged liquidation can succeed."

Actual staged liquidation can succeed. At G20 d327--336, \$49 cash, \$700 building salvage and \$665 mortgages pay a \$1,400 debt and leave \$14. Insolvency is also recognized correctly. At G16 d602, Gemini owes \$2,000 with \$288 and declares at the first debt menu without liquidating. Its note adds \$375 from the light-blue hotels and \$730 from mortgages to reach \$1,393, which equals the generous bound. An early research label had omitted collateral and was corrected. These contrasts separate local financing competence, genuine impossibility and analyst error. None evaluates the earlier choices that produced the debt.

%======================================================================================
\ReplacementBlock{FF-L02}{sections/finance\_inheritance.tex line 29}
% Replaces the sentence "The note also contains arithmetic errors."

The note opens by asserting that full liquidation cannot cover the debt, then recomputes and reaches the same five-sale route and \$43 remainder as the retained witness. It also miscounts the houses that would remain.

%======================================================================================
\ReplacementBlock{FF-B01}{sections/bargaining.tex lines 28 and 30}
% Replaces both paragraphs.

Author-only completion offers are rejected in 159/289 cases (55.0\%), compared with 38/140 (27.1\%) for recipient-only completion. About half of this pooled gap reflects where the offers sit in a thread. Among opening offers, rejection counts are 138/187 (73.8\%) versus 31/51 (60.8\%); among later offers, they are 21/102 (20.6\%) versus 7/89 (7.9\%). Author-only offers more often open a thread, where rejection is common in every category. The within-position differences are about thirteen percentage points, compared with twenty-eight points pooled. Offers completing neither group give a reference. They are rejected in 152/238 openings (63.9\%) and 27/170 later offers (15.9\%). Author-only offers sit ten and five points above that reference, and recipient-only offers three and eight points below it. Opening offers that would complete both groups are rejected in 34/48 cases (70.8\%), close to the author-only rate. These are conditional descriptions, not an identified effect of completion.

The contrast also varies by recipient. Pooled over thread position, Sol, Gemini and Grok reject author-only offers more often than recipient-only offers and accept them less often. Opus shows neither ordering. It rejects 3/52 against 6/42, accepts 7/52 against 5/42, and counters forty-two and thirty-one of them. Within position the cells are small and the orderings are not uniform. For Sol, Gemini and Grok the rejection ordering holds in five of six recipient-by-position cells and fails for Sol's openings (19/41 against 3/6). The acceptance ordering holds in three of six. It ties at zero for Sol's openings and reverses for Gemini's later offers (4/9 against 4/12) and for Grok's openings (3/61 against 0/18, where two of the three acceptances are text/call conflicts). The pooled rejection contrast holds in all five seed blocks and seventeen of nineteen games containing both categories; the acceptance contrast holds in four blocks and thirteen games.

%======================================================================================
\ReplacementBlock{FF-B02}{sections/bargaining.tex lines 56 and 58}
% Replaces the first paragraph of the subsection and adds one paragraph after the second.

Counter contents show how disagreement changes a transaction. Of 111 counters to author-only completion offers, ninety-two preserve the same original author's newly complete group, and the same ninety-two retain every deed that author requested. Eighty-eight still leave only that party completing. Of those eighty-eight, eighty-five improve the counterer's net cash, with a median increase of \$125, and seventy-nine preserve the entire noncash bundle, including deeds and jail cards. All seventy-nine increase the counterer's cash. These are cash repricings of the same exchange; the other nine also alter noncash terms.

% --- new paragraph after the paragraph ending "or equivalence between the populations." ---
Two configurations write most of these counters. Sol and Opus supply ninety-two of the 111 counters to author-only offers and eighty of the 110 in the comparison row, and Opus writes thirty-seven and twenty-three of them under action-only recovery. The form is still the majority for every counterer. An unchanged noncash bundle with more cash for the counterer accounts for 32/50 and 36/47 of Sol's counters in the two rows, 6/9 and 3/5 of Gemini's, 32/42 and 22/33 of Opus's, and 9/10 and 20/25 of Grok's. The 111 counters fall in fifty-nine threads and the 110 in sixty-four.

%======================================================================================
\ReplacementBlock{FF-B03}{sections/bargaining.tex line 76}
% Replaces the third through eighth sentences ("Grok and Gemini each supply ... across all four.").

Joining the two retained row ledgers itemizes the agreed sixty-one as twenty-nine of Gemini's thirty-one rejections, twenty-seven of Grok's thirty-one, all four of Sol's and one of Opus's eight. None of Opus's eight has a private note. Thirty-four of the sixty-one lie among the thirty-six offers that would also complete the author's group, where a rival monopoly is itself on offer, and twenty-seven among the thirty-eight recipient-only offers. The reasons also concern building funds or survival. The coding is fallible, and thirty-four of the seventy-four records come from two games. It documents Gemini and Grok, which produce most outright refusals, and says little about Opus.

%======================================================================================
\ReplacementBlock{FF-C01}{sections/finance\_inheritance.tex lines 88 and 89}
% Replaces the G18 finding and the second through fourth sentences of the reading.

\TranscriptFinding{No sale occurs at \$250. Six turns later Opus bids \$150 and then \$225 for the same deed, both below the withdrawn price. Sol, by then holding \$249, asks \$450, \$350 and \$300, and no sale follows. Opus redeems Pennsylvania two turns after the withdrawal and North Carolina two turns after that. It later receives Sol's estate and wins. The earlier assurance is present in both focal requests at invocation and withdrawal.}

\TranscriptReading{The record shows a repeated assurance, its invocation and changed terms after financial deterioration. The explanation cites the cost of prolonging a nearly bankrupt rival and of redeeming and developing the group. Opus then redeems two green deeds and bids again for Pacific within six turns, so the withdrawal was followed by a lower price and not by leaving the purchase. The explanation is a stated valuation, and the later conduct does not match its stated cost reason. No costly action in reliance or dishonesty when the original offer was made is identified. The engine supplies no binding standing-offer object, and the messages specify no expiry.}

%======================================================================================
\ReplacementBlock{FF-S01}{sections/setup.tex line 36}
% Insert after "An invalid action-only response receives no further corrective request."

The recovery policy also lists a fallback to a different model as a third stage. That stage never ran. All 1,696 Opus decisions return from the listed identifier, and none uses more than two attempts.

%======================================================================================
\ReplacementBlock{FF-S02}{sections/setup.tex line 40}
% Replaces the first sentence of the paragraph. The 1,068 comes from a simple content-name pattern
% and should be confirmed with the retained parser before use.

Among the 1,068 selected Grok responses whose content names an action, 139 name a different action from the received tool call.

%======================================================================================
\ReplacementBlock{FF-T01}{sections/temporal\_discussion.tex line 4; sections/methods.tex line 31}
% Replaces the first paragraph of the section.

The fixed sample holds 240 decisions. An earlier review drew three per player-game by hash rank from the decisions its public-message extractor had not flagged, so each configuration contributes sixty. It is not a random sample of all decisions. One LLM-assisted reader tagged fifty-nine sampled decisions as containing a forward-looking statement about the speaker's own action, thirty-six as containing a prospective statement about an opponent's capacity or outcome, and twenty-five as referring to correction or feedback, with overlap. Grok supplies thirty-nine of the fifty-nine and twenty-two of the thirty-six. Fifty-three private fields are absent, all for Opus, which prevents a comparable account of private planning across configurations.

% --- methods.tex line 31: replace "A single instance coded prospective statements in a fixed 240-decision sample.
%     Earlier review covered 614 overlapping catalog entries and 240 complete source bundles." with ---
A single instance coded prospective statements in a fixed sample of 240 decisions, three per player-game (Section~\ref{sec:temporal}).

%======================================================================================
\ReplacementBlock{FF-T02}{sections/temporal\_discussion.tex line 10}
% Replaces the third through fifth sentences ("Of fifty later submissions ... in the same turn.").

Of fifty later submissions after a decision containing a corrective request with the validator message, nine repeat it. Of fifty-nine after a decision without one, fifty-four repeat it, and all fifty-nine are Opus's, because the refusal path had already consumed the feedback opportunity. Feedback and configuration therefore cannot be separated in this record. States and targets also differ, and some successors occur in the same turn.

%======================================================================================
\ReplacementBlock{FF-R01}{sections/related\_limits.tex lines 7 and 9}
% Replaces the second through fifth sentences of the AgentMonopoly paragraph, and adds a locator in the next.

The paper's shared prompt encourages bluffing while prohibiting public messages that contradict the action taken, and the paper states that a player is eliminated when its cash becomes negative~\cite[\S4, App.~A.2]{wu2026agentmonopoly}. The repository implementation inspected at commit \texttt{2434f427} lets a trade recipient only accept or reject, has no debt-liquidation step before that elimination, and transfers buildings intact to a player-creditor~\cite{agentmonopoly2026repo}. Those three are implementation findings that the paper does not state.

% --- line 9: add the locator to the existing sentence ---
Its Monopoly return summaries use a bankruptcy treatment that sells buildings to the bank~\cite[App.~C.3]{hu2026finboardbench}, preventing direct comparison with receipts on inherited development here.

%======================================================================================
\ReplacementBlock{FF-W01}{manuscript.tex line 13 (abstract)}
% Replaces the third, fifth and tenth sentences of the abstract.

Sol, Gemini, Opus and Grok complete 17, 18, 15 and 13 color groups through trade, but build on 16, 17, 7 and 3 in their first management windows.

Opus receives action-only refusal recovery on 1,604 of its 1,696 decisions.

Eight contextual incidents connect these populations to rival-sensitive refusals, authorized bluffing, a selective disclosure, a financial-denial appeal and a changed assurance, alongside factual correction and contradictory execution.

%======================================================================================
\ReplacementBlock{FF-C02 (optional)}{sections/bargaining.tex lines 85 and 89}
% Adds one paragraph after the nine-house correction and extends the G2 finding.

The ask is specific to the buyer. One turn earlier Sol had countered Grok's \$200 bid for Marvin Gardens at \$450, and Grok, which held no yellow, answered that \$450 was the price for the player who completes yellow. Against Opus, Sol opens at \$650, comes down to \$600 and refuses \$450. The two negotiations differ in counterparty and order, so this is an observed contrast in one seller's asks, not a measured completion premium.

\TranscriptFinding{Sol accepts the added \$100 and completes yellow. Both deeds arrive mortgaged. Sol pays \$26 transfer interest and \$286 to redeem them, mortgages four other deeds for \$410, and spends \$300 on two yellow houses. Opus is eliminated three turns later, owing Gemini \$1,300, and its newly acquired railroads pass to Gemini with its developed light blues. On the next turn Sol lands on one of them, sells both houses for \$150 and mortgages two yellows to pay \$450. Sol collects no developed yellow rent and \$44 of undeveloped Atlantic rent at turn 77. Its later bankruptcy transfers the yellows to Gemini, the eventual winner, which redeems them and builds six houses the following turn.}

%======================================================================================
\ReplacementBlock{FF-C03 (optional)}{sections/bargaining.tex line 122}
% Extends the G7 finding with the ledger follow-through.

\TranscriptFinding{The deed transfers for \$600. Sol's note says it prefers reinvesting to holding out for another \$25. Grok completes the greens at turn 95 by winning a \$171 auction for Pennsylvania and never builds on them. At turn 110 it is the first player eliminated, owing Sol \$1,400, and North Carolina returns to Sol inside the complete group. Sol puts a house on each green two turns later.}

%======================================================================================
\ReplacementBlock{FF-C04 (optional)}{sections/finance\_inheritance.tex line 67}
% Insert after the first paragraph of Transcript 7.

The mortgage on Pennsylvania dates from Grok's first management window after completing yellow. At turn 63 Grok mortgaged Park Place and Pennsylvania for \$335, the content of its next response named a house on each of two yellows, and the received call ended the turn (Section~\ref{sec:development}). Grok builds nothing in this game.

%======================================================================================
\ReplacementBlock{FF-C05 (optional)}{sections/bargaining.tex line 134}
% Replaces the G19 finding.

\TranscriptFinding{The swap and \$150 payment execute and complete no group. Sol's note calls the exchange strictly favorable.}

%======================================================================================
\ReplacementBlock{FF-I01 (optional)}{sections/finance\_inheritance.tex lines 36 and 55}
% Adds a reference share after the first sentence of line 36 and a count in line 55.

One player in four wins each game, which gives a reference share of a quarter. Trade completions sit near it at 18/63 and auction completions at 2/9, while direct purchases (12/16) and nonterminal estates (45/52) sit far above it.

% --- line 55: replace "The creditor often already owns the developed property causing the debt" with ---
In fifty-seven of the sixty eliminations the triggering charge is rent on a developed color group that the creditor owns

%======================================================================================
\ReplacementBlock{FF-W02 (optional compressions)}{sections/bargaining.tex line 52; sections/development.tex line 34}

% bargaining.tex line 52: the four-sentence paragraph becomes
Low acceptance does not rule out agreement bias, because the record has no fair-price or utility reference against which a concession could be judged.

% development.tex line 34: final sentence becomes
Appendix~\ref{app:economy} gives the per-configuration receipts.
